\section{Power spectrum}
The H II power spectrum is the Fourier transform of the correlation function,
\[
P(k) = \int \dd^3{r} e^{i \textbf{k}\cdot \textbf{r}}\xi(r)
\]

Following the article, we discuss the power spectrum in two terms, 
\begin{itemize}
  \item Two-bubble term: the two contributions are from two bubbles;
  \item One-bubble term: the two contributions are from the same bubble.
\end{itemize}

\subsection{Two-bubble terms}
The two-bubble correlation function is
\[
  \xi^{2b}_{\delta_{HII}\delta_{HII}}
  =
  \xi^{2b}_{\delta x\delta x}(r)
  +x_e^2\xi^{2b}_{\delta\delta}(r)
  +\xi^{2b}_{\delta x \delta \delta x \delta}(r)
  +2 x_e\xi^{2b}_{\delta x\delta}(r),
\]
and the averaged pertubation is,
\[
\left<
\delta{x(\textbf{r})}
\right>_P
=
-\dfrac{
  (1-x_e)\ln(1-x_e)
  }{\left<
    V_b
  \right>} b \int\dd{R}{
    P(R) V_b(R) \delta_W(\mathbf{r})
  }.
\]
Note that $\delta_W(\mathbf{r})$ is given by
\[
\delta_W(\textbf{r}) = \int \dd^3{r^\prime}{\delta(\textbf{r}^\prime)W_R(\textbf{r}-\textbf{r}^\prime)},
\]
this is a convolution of matter fluctuation $\delta(\textbf{r})$ and window function $W_R(\textbf{r})$, which will convert to a product in the fourier space, that is,
\begin{align*}
  \delta_W(k;R) = \int \dd^3{r} e^{i \textbf{k}\cdot \textbf{r}}
  \left[
    \int \dd^3{r^\prime}{\delta(\textbf{r}^\prime)W_R(\textbf{r}-\textbf{r}^\prime)}
  \right]
  = \delta(k)\times W_R(k).
\end{align*}

So the Fourier transform of $\left<\delta{x(\textbf{r})}\right>_P$ is,
\[
\left<\delta{x(k)}\right> = 
-\dfrac{(1-x_e)\ln(1-x_e)}{\left<V_b\right>} 
b \int\dd{R}{
    P(R) V_b(R) \delta_W(k;R)
}
\equiv \dfrac{\alpha}{\left<V_b\right>}  \delta(k) 
\int\dd{R}{
    P(R) V_b(R) W_R(k)
}
= \dfrac{\alpha}\delta(k)\left<W_R(k)\right>,
\]
here $\alpha \equiv -(1-x_e)\ln(1-x_e)b$ and $\left<W_R(k)\right>\equiv \dfrac{\int\dd{R}{
    P(R) V_b(R) W_R(k)
}}{\left<V_b\right>}$.

Taking the Fourier transform of the corelation function, we obtain the power spectrum of the two-bubble term, and its component is given by
\begin{align*}
  P^{2b}_{\delta x\delta x}(k)
  = \int \dd^3{r} e^{i \mathbf{k}\cdot \mathbf{r}}\xi^{2b}_{\delta x\delta x}(r)
  = \int \dd^3{r} e^{i \mathbf{k}\cdot \mathbf{r}}
    \left<\delta x(\mathbf{r_1})\delta x(\mathbf{r_2})\right>
  = \int \dd^3{r} e^{i \mathbf{k}\cdot \mathbf{r}}
    \left<\delta x(\mathbf{r_1})\right>
    \left<\delta x(\mathbf{r_2})\right>
  = \alpha^2 
\end{align*}


\begin{align*}
  P^{2b}(k) = \left[(1-x_e)\ln(1-x_e)b\dfrac{\int\dd{R}P(R)V(R)W(kR)}{\int\dd{R}P(R)V(R)}-x_e\right]^2P(k)
\end{align*}
\begin{align*}
  F(k) = \dfrac{\int\dd{R}P(R)[V(R)]^2[W(kR)]^2}{\int\dd{R}P(R)V(R)}\\
  G(k) = \int\dfrac{\dd^3\mathbf{k}'}{(2\pi)^3}P(\left|\mathbf{k}-\mathbf{k}'\right|)F(k')\\
  P^{1b}(k) = x_e(1-x_e)[F(k)+G(k)]
\end{align*}


\begin{align*}
    P^{2b}_{\delta_{HII}}(k) = \left[(1-x_e)\ln(1-x_e)b\left<W_R(k)\right>-x_e\right]^2P(k)
\end{align*}
\begin{align*}
  P^{1b}_{\delta_{HII}}(k) &= x_e(1-x_e)[F(k)+G(k)]\\
  F(k) &= \left<V_b\right>\left<W_R^2(k)\right>\\
  G(k) &= \int\dfrac{\dd^3\mathbf{k}'}{(2\pi)^3}P(\left|\mathbf{k}-\mathbf{k}'\right|)F(k')
\end{align*}

\begin{align*}
  P_{\delta_{HII}}(k) = P^{2b}_{\delta_{HII}}(k)+P^{1b}_{\delta_{HII}}(k) 
\end{align*}

\begin{align*}
  x_e = \dfrac{1}{2}\left[1-\tanh\left(\dfrac{y(z)-y(z_{mid})}{\Delta y}\right)\right], y(z) = \left(1+z\right)^{\frac{3}{2}}
\end{align*}

In the 2-bubble power spectrum, the term 
\[
x_e^2 P(k)
\]
doesn't depend on the scale 


